\section{Benchmarking the two options}
\label{sec:setup}

\subsection{Fully replicated layout: near-identical CodeSearchNet collections}

\mypar{Dataset and system.} We use the
CodeSearchNet~\cite{husain2019codesearchnet} corpus from the VIBE
benchmark~\cite{vibe2024}, 768 dimensional Jina
embeddings~\cite{gunther2023jina} benchmarked with cosine. The
\textit{original} variant is the full corpus, and \textit{trimmed2} drops
the first 2\% of the corpus and recomputes ground truth against the
remaining vectors, so it behaves like a second tenant whose collection is
nearly but not byte-for-byte equal. Neighbor IDs from trimmed2 are
remapped into the original index space when we compare the two. The
system is Milvus~\cite{wang2021milvus} standalone built from source,
backed by etcd, MinIO, and rocksmq. Each experiment creates one or two
collections, inserts vectors in batches of 10{,}000, builds an
HNSW~\cite{malkov2018hnsw} index ($M{=}16$, $\text{efC}{=}200$), and
benchmarks Recall@100 at $\text{ef}{=}200$, following the evaluation
methodology of ANN-Benchmarks~\cite{aumuller2020annbenchmarks}.
Per-segment primary keys are read directly out of the MinIO Parquet
binlogs so we can compute exact vector-level overlap between any two
segments without modifying Milvus.

\mypar{Flush, jitter, and compaction in Milvus.} Milvus organizes every
collection as an append-only sequence of segments. Incoming vectors are
first buffered in a growing segment in memory; when that segment reaches
a target size it is \emph{flushed}, meaning it is sealed, written out to
object storage as Parquet binlogs, and subsequently becomes a candidate
for index building. Instead of sealing a growing segment at exactly the configured
threshold, Milvus applies a random \emph{jitter} so that each segment
is flushed around, but not exactly at, the target size. Jitter serves
two purposes: it desynchronizes flushes across data nodes so that many
parallel ingest streams do not hit the seal threshold at the same
moment and overwhelm the control plane, and it acts as a tie breaker
when several seal policies fire together so that not every segment is
sealed at an identical size. The side effect on a single-node
standalone run is that two otherwise identical ingest streams produce
segments whose vector counts, and therefore whose vector-range boundaries,
differ from run to run. Once segments are sealed, a background
\emph{compactor} periodically inspects them and may merge several small
or partially deleted segments into a new, larger one, rewriting their
binlogs in the process. Compaction keeps the number of segments bounded
and reclaims space from deleted vectors, but from the point of view of
cross-collection comparison it is another source of shuffling: a vector
that was originally flushed into segment $A$ may end up physically
stored in a compacted segment $A'$ whose boundaries are unrelated to
$A$'s original vector range.

\mypar{Configurations.} We test three variants of Milvus that progressively
remove non-determinism from the ingest path.
\textbf{bothOn} is stock Milvus, segment size at flush time is jittered
around the target and a background compactor may merge and rewrite small
segments. \textbf{nojitter} disables the random jitter, every flush
targets the exact same segment size, but background compaction still
runs. \textbf{noJitterCompaction} additionally disables compaction, so
each segment is exactly one flush worth of vectors and boundaries are
fully determined by insert order and flush size.

\subsection{Shared partition layout: split the  Wikipedia workload}

\begin{table}[t]
\centering
\caption{Three workload variants for the shared partition layout.
$\mathrm{QC}$ is the fraction of queries whose top 100 ground truth
neighbors are predominantly served from the shared partition, and
the scatter ratio gives the share of cold vectors assigned to the
shared, $A$, and $B$ partitions.}
\label{tab:wiki-variants}
\small
\setlength{\tabcolsep}{4pt}
\begin{tabular}{lccrrr}
\toprule
variant & QC & scatter & shared & A & B \\
\midrule
QC02\_SS02 & 0.2 & 0.2/0.4/0.4 & 7.1M  & 14.2M & 14.2M \\
QC02\_SS08 & 0.2 & 0.8/0.1/0.1 & 28.0M & 3.9M  & 3.9M  \\
QC08\_SS08 & 0.8 & 0.8/0.1/0.1 & 28.1M & 3.7M  & 3.7M  \\
\bottomrule
\end{tabular}
\end{table}

\begin{table}[t]
\centering
\caption{The six loaded collections, one per workload variant
under each of the two compaction settings. All collections are
loaded via mmap. The last column lists segment counts in the
shared, $A$, and $B$ partitions.}
\label{tab:wiki-collections}
\small
\setlength{\tabcolsep}{4pt}
\begin{tabular}{lrr}
\toprule
collection & segs & shared / A / B \\
\midrule
QC02\_SS02            & 275 & 55  / 111 / 109 \\
QC02\_SS02\_nocompact & 861 & 171 / 345 / 345 \\
QC02\_SS08            & 271 & 209 / 31  / 31  \\
QC02\_SS08\_nocompact & 868 & 678 / 95  / 95  \\
QC08\_SS08            & 264 & 204 / 29  / 31  \\
QC08\_SS08\_nocompact & 860 & 680 / 90  / 90  \\
\bottomrule
\end{tabular}
\end{table}

\mypar{Pipeline.} We use a corpus of 35M Wikipedia chunks embedded
with the Cohere v3 model and reembedded with BAAI
\texttt{bge-base-en-v1.5} at 768 dimensions, together with 15{,}000
queries drawn from the TriviaQA validation set. Ground truth is
computed by exact $k$ nearest neighbor search at $k{=}100$ against
the full corpus. The vector index is HNSW~\cite{malkov2018hnsw} with
$M{=}16$, $\text{efC}{=}200$, and cosine distance, matching the
configuration used in the fully replicated experiments.

\mypar{Building the splits.} Each workload variant splits the
corpus into one shared partition and two tenant partitions $A$ and
$B$. The shared partition is grown greedily from the most frequently
accessed ground truth vectors until a target fraction $\mathrm{QC}$
of the queries has at least 50 of its top 100 ground truth neighbors
inside the shared partition, that is, until $\mathrm{QC}$ of the
queries are predominantly answerable from the shared corpus. The
queries are then split evenly between tenants $A$ and $B$, and each
query's remaining top 100 neighbors are routed into its own tenant
partition. Cold vectors that no query touches in its top 100 are
distributed across the shared partition, $A$, and $B$ according to a
fixed scatter ratio.

\mypar{Three workload variants.} We instantiate three variants that
differ in $\mathrm{QC}$ and in the cold-vector scatter ratio, listed
in \autoref{tab:wiki-variants}. \textbf{QC02\_SS02} keeps a small
shared partition with relatively even cold-vector scatter, so each
tenant partition holds 14.2M vectors. \textbf{QC02\_SS08} retains the
same query coverage but pushes most cold vectors into the shared
partition, producing a 28.0M vector shared partition and 3.9M vector
tenant partitions. \textbf{QC08\_SS08} additionally raises query
coverage so that the bulk of queries can be answered from the
shared partition while the tenant partitions stay small.

\mypar{Loading modes and compaction.} All variants are loaded with
vectors and HNSW index files mmapped so that the operating system
pages data in on demand rather than holding the whole collection in
heap. We do not also evaluate a fully in-heap mode because the
QC02\_SS08 and QC08\_SS08 variants carry shared partitions of
roughly 28 million vectors that already exceed the RAM cap of the
test machine, so an in-heap configuration is not feasible on those
variants and a heap-versus-mmap comparison would only be possible
on the smallest variant. We therefore vary only the compaction
axis: each variant is loaded once with compaction enabled at the
Milvus default and once with compaction disabled, leaving every
segment as a single flush worth of vectors. The result is the six
collections summarized in \autoref{tab:wiki-collections}.
Compaction-disabled collections produce roughly three times as many
segments as their compacted counterparts at the same total vector
count, for example 861 segments versus 275 segments on QC02\_SS02.

\mypar{Milvus configuration.} The Milvus server is configured with
mmap enabled for both vectors and index files at a 512\,MB threshold, with the maximum segment size and a sealed ratio
of 0.24, and a 96\,GB RAM cap. These settings keep the
working set close to but not strictly inside main memory, so any
cross-tenant resource pressure surfaces in the throughput
measurements rather than being hidden by overprovisioned RAM.
